% Options for packages loaded elsewhere
% Options for packages loaded elsewhere
\PassOptionsToPackage{unicode}{hyperref}
\PassOptionsToPackage{hyphens}{url}
\PassOptionsToPackage{dvipsnames,svgnames,x11names}{xcolor}
%
\documentclass[
  letterpaper,
  DIV=11,
  numbers=noendperiod]{scrreprt}
\usepackage{xcolor}
\usepackage{amsmath,amssymb}
\setcounter{secnumdepth}{5}
\usepackage{iftex}
\ifPDFTeX
  \usepackage[T1]{fontenc}
  \usepackage[utf8]{inputenc}
  \usepackage{textcomp} % provide euro and other symbols
\else % if luatex or xetex
  \usepackage{unicode-math} % this also loads fontspec
  \defaultfontfeatures{Scale=MatchLowercase}
  \defaultfontfeatures[\rmfamily]{Ligatures=TeX,Scale=1}
\fi
\usepackage{lmodern}
\ifPDFTeX\else
  % xetex/luatex font selection
\fi
% Use upquote if available, for straight quotes in verbatim environments
\IfFileExists{upquote.sty}{\usepackage{upquote}}{}
\IfFileExists{microtype.sty}{% use microtype if available
  \usepackage[]{microtype}
  \UseMicrotypeSet[protrusion]{basicmath} % disable protrusion for tt fonts
}{}
\makeatletter
\@ifundefined{KOMAClassName}{% if non-KOMA class
  \IfFileExists{parskip.sty}{%
    \usepackage{parskip}
  }{% else
    \setlength{\parindent}{0pt}
    \setlength{\parskip}{6pt plus 2pt minus 1pt}}
}{% if KOMA class
  \KOMAoptions{parskip=half}}
\makeatother
% Make \paragraph and \subparagraph free-standing
\makeatletter
\ifx\paragraph\undefined\else
  \let\oldparagraph\paragraph
  \renewcommand{\paragraph}{
    \@ifstar
      \xxxParagraphStar
      \xxxParagraphNoStar
  }
  \newcommand{\xxxParagraphStar}[1]{\oldparagraph*{#1}\mbox{}}
  \newcommand{\xxxParagraphNoStar}[1]{\oldparagraph{#1}\mbox{}}
\fi
\ifx\subparagraph\undefined\else
  \let\oldsubparagraph\subparagraph
  \renewcommand{\subparagraph}{
    \@ifstar
      \xxxSubParagraphStar
      \xxxSubParagraphNoStar
  }
  \newcommand{\xxxSubParagraphStar}[1]{\oldsubparagraph*{#1}\mbox{}}
  \newcommand{\xxxSubParagraphNoStar}[1]{\oldsubparagraph{#1}\mbox{}}
\fi
\makeatother

\usepackage{color}
\usepackage{fancyvrb}
\newcommand{\VerbBar}{|}
\newcommand{\VERB}{\Verb[commandchars=\\\{\}]}
\DefineVerbatimEnvironment{Highlighting}{Verbatim}{commandchars=\\\{\}}
% Add ',fontsize=\small' for more characters per line
\usepackage{framed}
\definecolor{shadecolor}{RGB}{241,243,245}
\newenvironment{Shaded}{\begin{snugshade}}{\end{snugshade}}
\newcommand{\AlertTok}[1]{\textcolor[rgb]{0.68,0.00,0.00}{#1}}
\newcommand{\AnnotationTok}[1]{\textcolor[rgb]{0.37,0.37,0.37}{#1}}
\newcommand{\AttributeTok}[1]{\textcolor[rgb]{0.40,0.45,0.13}{#1}}
\newcommand{\BaseNTok}[1]{\textcolor[rgb]{0.68,0.00,0.00}{#1}}
\newcommand{\BuiltInTok}[1]{\textcolor[rgb]{0.00,0.23,0.31}{#1}}
\newcommand{\CharTok}[1]{\textcolor[rgb]{0.13,0.47,0.30}{#1}}
\newcommand{\CommentTok}[1]{\textcolor[rgb]{0.37,0.37,0.37}{#1}}
\newcommand{\CommentVarTok}[1]{\textcolor[rgb]{0.37,0.37,0.37}{\textit{#1}}}
\newcommand{\ConstantTok}[1]{\textcolor[rgb]{0.56,0.35,0.01}{#1}}
\newcommand{\ControlFlowTok}[1]{\textcolor[rgb]{0.00,0.23,0.31}{\textbf{#1}}}
\newcommand{\DataTypeTok}[1]{\textcolor[rgb]{0.68,0.00,0.00}{#1}}
\newcommand{\DecValTok}[1]{\textcolor[rgb]{0.68,0.00,0.00}{#1}}
\newcommand{\DocumentationTok}[1]{\textcolor[rgb]{0.37,0.37,0.37}{\textit{#1}}}
\newcommand{\ErrorTok}[1]{\textcolor[rgb]{0.68,0.00,0.00}{#1}}
\newcommand{\ExtensionTok}[1]{\textcolor[rgb]{0.00,0.23,0.31}{#1}}
\newcommand{\FloatTok}[1]{\textcolor[rgb]{0.68,0.00,0.00}{#1}}
\newcommand{\FunctionTok}[1]{\textcolor[rgb]{0.28,0.35,0.67}{#1}}
\newcommand{\ImportTok}[1]{\textcolor[rgb]{0.00,0.46,0.62}{#1}}
\newcommand{\InformationTok}[1]{\textcolor[rgb]{0.37,0.37,0.37}{#1}}
\newcommand{\KeywordTok}[1]{\textcolor[rgb]{0.00,0.23,0.31}{\textbf{#1}}}
\newcommand{\NormalTok}[1]{\textcolor[rgb]{0.00,0.23,0.31}{#1}}
\newcommand{\OperatorTok}[1]{\textcolor[rgb]{0.37,0.37,0.37}{#1}}
\newcommand{\OtherTok}[1]{\textcolor[rgb]{0.00,0.23,0.31}{#1}}
\newcommand{\PreprocessorTok}[1]{\textcolor[rgb]{0.68,0.00,0.00}{#1}}
\newcommand{\RegionMarkerTok}[1]{\textcolor[rgb]{0.00,0.23,0.31}{#1}}
\newcommand{\SpecialCharTok}[1]{\textcolor[rgb]{0.37,0.37,0.37}{#1}}
\newcommand{\SpecialStringTok}[1]{\textcolor[rgb]{0.13,0.47,0.30}{#1}}
\newcommand{\StringTok}[1]{\textcolor[rgb]{0.13,0.47,0.30}{#1}}
\newcommand{\VariableTok}[1]{\textcolor[rgb]{0.07,0.07,0.07}{#1}}
\newcommand{\VerbatimStringTok}[1]{\textcolor[rgb]{0.13,0.47,0.30}{#1}}
\newcommand{\WarningTok}[1]{\textcolor[rgb]{0.37,0.37,0.37}{\textit{#1}}}

\usepackage{longtable,booktabs,array}
\newcounter{none} % for unnumbered tables
\usepackage{calc} % for calculating minipage widths
% Correct order of tables after \paragraph or \subparagraph
\usepackage{etoolbox}
\makeatletter
\patchcmd\longtable{\par}{\if@noskipsec\mbox{}\fi\par}{}{}
\makeatother
% Allow footnotes in longtable head/foot
\IfFileExists{footnotehyper.sty}{\usepackage{footnotehyper}}{\usepackage{footnote}}
\makesavenoteenv{longtable}
\usepackage{graphicx}
\makeatletter
\newsavebox\pandoc@box
\newcommand*\pandocbounded[1]{% scales image to fit in text height/width
  \sbox\pandoc@box{#1}%
  \Gscale@div\@tempa{\textheight}{\dimexpr\ht\pandoc@box+\dp\pandoc@box\relax}%
  \Gscale@div\@tempb{\linewidth}{\wd\pandoc@box}%
  \ifdim\@tempb\p@<\@tempa\p@\let\@tempa\@tempb\fi% select the smaller of both
  \ifdim\@tempa\p@<\p@\scalebox{\@tempa}{\usebox\pandoc@box}%
  \else\usebox{\pandoc@box}%
  \fi%
}
% Set default figure placement to htbp
\def\fps@figure{htbp}
\makeatother

\ifLuaTeX
  \usepackage{luacolor}
  \usepackage[soul]{lua-ul}
\else
  \usepackage{soul}
\fi




\setlength{\emergencystretch}{3em} % prevent overfull lines

\providecommand{\tightlist}{%
  \setlength{\itemsep}{0pt}\setlength{\parskip}{0pt}}



 


\KOMAoption{captions}{tableheading}
\makeatletter
\@ifpackageloaded{bookmark}{}{\usepackage{bookmark}}
\makeatother
\makeatletter
\@ifpackageloaded{caption}{}{\usepackage{caption}}
\AtBeginDocument{%
\ifdefined\contentsname
  \renewcommand*\contentsname{Table of contents}
\else
  \newcommand\contentsname{Table of contents}
\fi
\ifdefined\listfigurename
  \renewcommand*\listfigurename{List of Figures}
\else
  \newcommand\listfigurename{List of Figures}
\fi
\ifdefined\listtablename
  \renewcommand*\listtablename{List of Tables}
\else
  \newcommand\listtablename{List of Tables}
\fi
\ifdefined\figurename
  \renewcommand*\figurename{Figure}
\else
  \newcommand\figurename{Figure}
\fi
\ifdefined\tablename
  \renewcommand*\tablename{Table}
\else
  \newcommand\tablename{Table}
\fi
}
\@ifpackageloaded{float}{}{\usepackage{float}}
\floatstyle{ruled}
\@ifundefined{c@chapter}{\newfloat{codelisting}{h}{lop}}{\newfloat{codelisting}{h}{lop}[chapter]}
\floatname{codelisting}{Listing}
\newcommand*\listoflistings{\listof{codelisting}{List of Listings}}
\makeatother
\makeatletter
\makeatother
\makeatletter
\@ifpackageloaded{caption}{}{\usepackage{caption}}
\@ifpackageloaded{subcaption}{}{\usepackage{subcaption}}
\makeatother
\usepackage{bookmark}
\IfFileExists{xurl.sty}{\usepackage{xurl}}{} % add URL line breaks if available
\urlstyle{same}
% fallback for those not using the hyperref driver hyperxmp:
\makeatletter
\@ifundefined{xmpquote}{\newcommand{\xmpquote}[1]{#1}}{}
\makeatother
\hypersetup{
  pdftitle={quantcrim},
  pdfauthor={Joost Jansen; Gijs Custers},
  colorlinks=true,
  linkcolor={blue},
  filecolor={Maroon},
  citecolor={Blue},
  urlcolor={Blue},
  pdfcreator={LaTeX via pandoc}}


\title{quantcrim}
\usepackage{etoolbox}
\makeatletter
\providecommand{\subtitle}[1]{% add subtitle to \maketitle
  \apptocmd{\@title}{\par {\large #1 \par}}{}{}
}
\makeatother
\subtitle{Een open-web handboek voor kwantitatief criminologisch
onderzoek}
\author{Joost Jansen \and Gijs Custers}
\date{2026-09-24}
\begin{document}
\maketitle

\renewcommand*\contentsname{Table of contents}
{
\setcounter{tocdepth}{2}
\tableofcontents
}

\bookmarksetup{startatroot}

\chapter{Voorwoord}\label{voorwoord}

\bookmarksetup{startatroot}

\chapter*{Voorwoord}\label{voorwoord-1}
\addcontentsline{toc}{chapter}{Voorwoord}

\markboth{Voorwoord}{Voorwoord}

Welkom bij het \href{quantcrim.pdf}{quantcrim handboek} Criminologie.
Hier ga je leren hoe je kwantitatieve data kan analyseren met R. Het
boek is ontworpen en beschikbaar gesteld door docenten en onderzoekers
van de opleiding Criminologie aan de Erasmus Universiteit. Het boek is
bedoeld voor zowel bachelor- als masterstudenten.

Hieronder vind je een voorbeeld van hoe makkelijk het is om data in te
laden, een regressiemodel te schatten en de resultaten ervan te
presenteren in een plot.\footnote{We gaan in een later hoofdstuk zien
  wat een regressiemodel precies is.}

\begin{Shaded}
\begin{Highlighting}[]
\CommentTok{\# test}
\end{Highlighting}
\end{Shaded}

\section{Voor we verder gaan..}\label{voor-we-verder-gaan..}

We hopen dat je dit boek handig vindt om mee te werken. In het kader van
Open Science hebben we geprobeerd dit boek volledig transparant en
reproduceerbaar te ontwikkelen. De broncode is beschikbaar via
\href{https://github.com/joostjansen/QuantCrim}{https://github.com/joostjansen/quantcrim}.
Als je vragen, suggesties of aanpassingen hebt, aarzel niet deze te
stellen.\footnote{Het beste is om dit via GitHub te doen in onze
  Repository.}

\bookmarksetup{startatroot}

\chapter{Voorbereiding}\label{voorbereiding}

\bookmarksetup{startatroot}

\chapter*{Voorbereidingen}\label{voorbereidingen}
\addcontentsline{toc}{chapter}{Voorbereidingen}

\markboth{Voorbereidingen}{Voorbereidingen}

Voor we tijdens de lessen aan de slag kunnen met het verrichten van
kwantitatief onderzoek, moet de vereiste software geinstalleerd worden.
We gaan werken met R (spreek uit op z'n Engels). R is een
programmeertaal die veel statistici gebruiken om hun analyses mee te
doen. De omgeving waarin je met R kan werken, heet RStudio. Zowel R als
RStudio moeten apart geinstalleerd worden. Je vindt de installatielinks
hieronder:

\begin{enumerate}
\def\labelenumi{\arabic{enumi}.}
\tightlist
\item
  \href{https://cran.r-project.org/}{R}
\item
  \href{https://docs.posit.co/ide/user/\#rstudio-ide-oss-downloads}{RStudio}
\end{enumerate}

Na installatie van de benodigde software moeten we R gereed maken door
de juiste packages (ook wel: libraries) te installeren. Veel
databewerkingen en analyses vergen namelijk speciaal daarvoor software.
Een package waarmee je mooie visualisaties kan maken van je data heet
bijvoorbeeld ggplot. Je installeert zo'n package door RStudio te openen
en in je console de volgende code te typen en dan te `runnen':

\begin{Shaded}
\begin{Highlighting}[]
\CommentTok{\# install.packages("ggplot2")}
\end{Highlighting}
\end{Shaded}

Als je het package geinstalleerd hebt (eenmalig) laad je in het vervolg
aan het begin van elke sessie het package in. Let op: in tegenstelling
tot bij het installeren van een package, gebruik je bij het laden ervan
geen aanhalingstekens. R is zeer gevoelig voor het verkeerd gebruik van
leestekens. Wees daar dus scherp op.

\begin{Shaded}
\begin{Highlighting}[]
\FunctionTok{library}\NormalTok{(ggplot2)}
\end{Highlighting}
\end{Shaded}

Hieronder vind je een lijst met alle packages die je tijdens dit vak
nodig hebt.

\begin{Shaded}
\begin{Highlighting}[]
\CommentTok{\# }
\CommentTok{\# install.packages()}
\CommentTok{\# install.packages()}
\end{Highlighting}
\end{Shaded}

\part{Data beschrijven}

\chapter{Data}\label{data}

\section{\texorpdfstring{\textbf{Wat zijn
data?}}{Wat zijn data?}}\label{wat-zijn-data}

In de criminologie zijn we geïnteresseerd in diverse vragen die we
kunnen beantwoorden aan de hand van statistisch onderzoek. Denk aan de
volgende vragen:

\begin{itemize}
\tightlist
\item
  Hoeveel inbraken vinden er jaarlijks plaats in Kralingen?
\item
  Wat is het gemiddelde geldbedrag waarvoor Nederlandse banken hebben
  geschikt in witwaszaken in 2025?
\item
  Hoeveel misdaadjournalisten zijn er vermoord het afgelopen decennium?
\item
  Wat is de associatie tussen drugsverslaving en delinquentie?
\item
  In hoeverre leidt seksuele straatintimidatie tot het vermijden van
  bepaalde plekken?
\end{itemize}

Om deze vragen te kunnen beantwoorden, werken we met \textbf{datasets}.
Een set is in de wiskunde een verzameling van dingen. Een dataset is dus
een verzameling van data. Het woord \textbf{data} is afkomstig uit het
Latijnse `datum' wat letterlijk `gegeven' betekent. Een dataset is dus
een verzameling gegevens. Dit kan gaan over gegevens van mensen
(bijvoorbeeld hun lengte of geslacht), maar bijvoorbeeld ook over andere
dingen zoals organisaties, steden of landen. Hetgene waar de data
informatie over verschaffen (mensen, organisaties, steden, landen)
noemen we de \textbf{analyse-eenheid}.

In deze cursus werken we steeds met dezelfde dataset. Dit is gebaseerd
op een onderzoek onder uitgaande jongeren, gebaseerd op het
\href{https://www.trimbos.nl/kennisbank/rapport-het-grote-uitgaansonderzoek-2023/}{Het
Grote Uitgaansonderzoek} van het \href{https://www.trimbos.nl/}{Trimbos
Instituut}. In dit onderzoek wordt onder andere gekeken hoe vaak
jongeren uitgaan, naar wat soort feesten en gelegenheden, wat voor
middelen ze gebruiken (bijvoorbeeld alcohol of andere drugs) en hoe ze
denken over drugsgebruik. Om de data te verzamelen is een online enquête
afgenomen via een gerichte wervingscampagne.

Laten we eerst kijken hoe zo'n dataset eruitziet. Dat doen we met behulp
van R. Zoals aangegeven in de introductie van dit online boek, het is
niet erg als je nog niet snapt hoe dit werkt. Daarvoor volg je uiteraard
deze cursus :)

\begin{Shaded}
\begin{Highlighting}[]
\CommentTok{\# laden packages}
\FunctionTok{library}\NormalTok{(readr)}
\FunctionTok{library}\NormalTok{(gt)}

\CommentTok{\# laden van de data}
\NormalTok{data }\OtherTok{\textless{}{-}} \FunctionTok{read\_rds}\NormalTok{(}\StringTok{"data/middelengebruik\_uitgaanders.rds"}\NormalTok{)}

\CommentTok{\# data inspecteren}
\FunctionTok{head}\NormalTok{(data[}\DecValTok{1}\SpecialCharTok{:}\DecValTok{5}\NormalTok{,}\DecValTok{2}\SpecialCharTok{:}\DecValTok{6}\NormalTok{]) }
\end{Highlighting}
\end{Shaded}

\begin{verbatim}
  geslacht leeftijd opleidingsniveau alcohol_ooit alcohol_laatste_jaar
1      man       22             hoog            1                    1
2    vrouw       23             hoog            1                    1
3    vrouw       18           midden            1                    0
4    vrouw       28             hoog            1                    1
5      man       19             hoog            1                    1
\end{verbatim}

\begin{Shaded}
\begin{Highlighting}[]
\CommentTok{\#|\textgreater{} \# eerste 5 rijen en kolommen}
\CommentTok{\#gt() |\textgreater{}}
  \CommentTok{\#tab\_header(title = "Voorbeeld van de data") |\textgreater{} \# titel geven aan output}
  \CommentTok{\#opt\_stylize(style = 6, color = "blue") \# lay{-}out bepalen}
\end{Highlighting}
\end{Shaded}

Hierboven staat een `sneak preview' van de dataset. Waar kijken we nu
eigenlijk naar? De dataset verschaft informatie die we verzameld hebben
over jongeren die hebben meegedaan aan de enquête. Elke rij
representeert een respondent: een jongere die de enquête heeft ingevuld.
De rijen zijn onze \textbf{observaties}. Elke kolom bevat de waarde van
die respondent met betrekking tot een kenmerk, bijvoorbeeld de waarde
die hoort bij een antwoord dat een respondent gaf op een in de enquête
gestelde vraag. De waarden van respondenten op kenmerken kunnen variëren
van respondent tot respondent. Niet iedereen is even oud, gebruikt even
vaak drugs, heeft dezelfde opleiding etc. De variatie in een bepaald
kenmerk meten we met behulp van \textbf{variabelen}. Deze zien we terug
in de kolommen. De kern van statistiek is dat we alle variatie in
kenmerken gaan gebruiken om vragen zoals die hierboven te kunnen
beantwoorden.

We kunnen bijvoorbeeld zien dat de eerste respondent een xxx is, xx jaar
oud, en een xxx opleidingsniveau heeft.

De vierde en vijfde kolom geven aan of iemand ooit alcohol heeft
gedronken of in het afgelopen jaar. Hier staat een `1' voor ja en een
`0' voor nee. Dit zul je vaker tegenkomen in het werken met datasets,
dat cijfers worden gebruikt om kwalitatieve informatie uit te drukken.
Hier komen we later nog meer op terug. De voorbeeld data toont dat vier
respondenten het afgelopen jaar alcohol hebben gedronken en één
respondent niet.

\chapter{Meetniveaus en
visualisaties}\label{meetniveaus-en-visualisaties}

Op de vorige pagina heb je gezien wat een dataset is, en dat deze
bestaat uit observaties (de rijen) en variabelen (de kolommen). In de
basis zijn er vier manieren met variabelen eigenschappen meetbaar te
maken (ook wel: operationaliseren). We noemen deze manieren
\textbf{meetniveaus.} Niet ieder meetniveau geeft even veel informatie.
Als we informatie willen over de leeftijden van onze respondenten,
kunnen we dit op verschillende manieren uitvragen. We kunnen direct
vragen naar hun leeftijd, maar we kunnen ze ook indelen in
leeftijdscategorieen (bijv.: jonger dan 18 en 18 jaar en ouder). We
bespreken hier de vier meetniveaus en daarbij gaan we meteen in op hoe
we de informatie die de variabelen verschaffen het beste visueel kunnen
presenteren.

Doorgaans maken we onderscheid tussen \textbf{kwalitatieve} en
\textbf{kwantitatieve} variabelen.

\section{Kwalitatieve variabelen}\label{kwalitatieve-variabelen}

Kwalitatieve variabelen bestaan uit categorieën die een kwalitatieve
eigenschap van een observatie beschrijven. Deze eigenschappen zijn niet
direct in cijfers uit te drukken, zoals bij kwantitatieve variabelen.
Een belangrijk aspect is dat de categorieën \emph{discreet} zijn: ze
moeten elkaar uitsluiten.

\subsection{Staafdiagram}\label{staafdiagram}

Laten we naar een voorbeeld kijken: geslacht. Stel dat we willen weten
hoeveel mannen, vrouwen of anders-identificerenden er in de dataset
zitten. Daarvoor kunnen we \textbf{staafdiagram} maken.

\begin{Shaded}
\begin{Highlighting}[]
\CommentTok{\# laden packages}
\FunctionTok{library}\NormalTok{(readr)}
\FunctionTok{library}\NormalTok{(ggplot2)}
\FunctionTok{library}\NormalTok{(tidyverse)}

\CommentTok{\# laden van de data}
\NormalTok{data }\OtherTok{\textless{}{-}} \FunctionTok{read\_rds}\NormalTok{(}\StringTok{"data/middelengebruik\_uitgaanders.rds"}\NormalTok{)}

\CommentTok{\# staafdiagram maken}
\FunctionTok{ggplot}\NormalTok{(data, }\FunctionTok{aes}\NormalTok{(}\AttributeTok{x =}\NormalTok{ geslacht)) }\SpecialCharTok{+} \CommentTok{\# data en variabele selecteren}
  \FunctionTok{geom\_bar}\NormalTok{() }\CommentTok{\# staafdiagram}
\end{Highlighting}
\end{Shaded}

\pandocbounded{\includegraphics[keepaspectratio]{002-02-meetniveaus-visualisaties_files/figure-pdf/unnamed-chunk-1-1.pdf}}

De staafdiagram geeft een visuele presentatie van het aantal mannen,
vrouwen en anders-identificerenden in de dataset. Het mooie aan R is dat
we de figuur visueel nog wat aantrekkelijker kunnen maken. Hiermee ga je
meer oefenen tijdens de cursus.

\begin{Shaded}
\begin{Highlighting}[]
\CommentTok{\# staafdiagram met betere visualisatie}
\FunctionTok{ggplot}\NormalTok{(data, }\FunctionTok{aes}\NormalTok{(}\AttributeTok{x =}\NormalTok{ geslacht)) }\SpecialCharTok{+}
  \FunctionTok{geom\_bar}\NormalTok{() }\SpecialCharTok{+}
  \FunctionTok{theme\_bw}\NormalTok{() }\SpecialCharTok{+} \CommentTok{\# witte achtergrond}
  \FunctionTok{ylab}\NormalTok{(}\StringTok{"aantal"}\NormalTok{)  }\CommentTok{\# label voor y{-}as}
\end{Highlighting}
\end{Shaded}

\pandocbounded{\includegraphics[keepaspectratio]{002-02-meetniveaus-visualisaties_files/figure-pdf/unnamed-chunk-2-1.pdf}}

De figuur heeft nu een witte achtergrond en de y-as heeft een
Nederlandstalig label.

\subsection{Simpele tabel}\label{simpele-tabel}

Als je wilt weten hoeveel respondenten er precies in elke categorie
zitten, kun je ook een simpele tabel maken.

\begin{Shaded}
\begin{Highlighting}[]
\CommentTok{\# tabel van geslacht}
\FunctionTok{table}\NormalTok{(data}\SpecialCharTok{$}\NormalTok{geslacht)}
\end{Highlighting}
\end{Shaded}

\begin{verbatim}

   man  vrouw anders 
   178    170      2 
\end{verbatim}

De dataset bevat in totaal dus xxx mannen, xxx vrouwen en xxx in de
categorie ``anders''.

\subsection{Nominale variabele}\label{nominale-variabele}

Geslacht is een voorbeeld van een \textbf{nominale} variabele. Dit is
een kwalitatieve variabele waarin geen rangorde bestaat tussen de
categorieën. De variabele bestaat uit kwalitatieve eigenschappen die van
elkaar verschillen, maar je kunt niet zeggen dat de ene categorie
intrinsiek `hoger' of `lager' is dan de ander.

\subsection{Taartdiagram}\label{taartdiagram}

Een andere manier om een nominale variabele te presenteren, is via een
\textbf{taartdiagram}. Een taartdiagram kan snel een overzicht geven van
de verhouding tussen categorieën binnen een bepaalde variabele.
Doorgaans werkt een taartdiagram het beste als er een beperkt aantal
categorieën zijn.

Hieronder staat een voorbeeld van taartdiagram van de variabele
\emph{manier\_verkrijgen\_drugs}. Dit gaat over de wijze waarop
uitgaande jongeren meestal aan hun drugs komen.

\begin{Shaded}
\begin{Highlighting}[]
\CommentTok{\# samenvatten van de relevante data}
\NormalTok{drugs\_verkrijgen }\OtherTok{\textless{}{-}}\NormalTok{ data }\SpecialCharTok{|\textgreater{}}
  \FunctionTok{filter}\NormalTok{(}\SpecialCharTok{!}\FunctionTok{is.na}\NormalTok{(manier\_verkrijgen\_drugs)) }\SpecialCharTok{|\textgreater{}} \CommentTok{\# missende waarden verwijderen}
  \FunctionTok{group\_by}\NormalTok{(manier\_verkrijgen\_drugs) }\SpecialCharTok{|\textgreater{}} \CommentTok{\# groeperen op variabele}
  \FunctionTok{summarise}\NormalTok{(}\AttributeTok{n =} \FunctionTok{n}\NormalTok{()) }\SpecialCharTok{|\textgreater{}} \CommentTok{\# aantal respondenten per categorie berekenen}
  \FunctionTok{mutate}\NormalTok{(}\AttributeTok{perc =} \FunctionTok{round}\NormalTok{(n}\SpecialCharTok{/}\FunctionTok{sum}\NormalTok{(n) }\SpecialCharTok{*} \DecValTok{100}\NormalTok{, }\DecValTok{1}\NormalTok{)) }\CommentTok{\# percentages berekenen en afronden op 1 decimaal}

\CommentTok{\# taartdiagram van hoe drugs worden verkregen}
\FunctionTok{ggplot}\NormalTok{(drugs\_verkrijgen, }\FunctionTok{aes}\NormalTok{(}\AttributeTok{x=} \StringTok{""}\NormalTok{, }\AttributeTok{y =}\NormalTok{ perc, }\AttributeTok{fill =}\NormalTok{ manier\_verkrijgen\_drugs)) }\SpecialCharTok{+} 
  \CommentTok{\# percentages (y) en categorieen (fill) selecteren}
  \FunctionTok{geom\_bar}\NormalTok{(}\AttributeTok{stat =} \StringTok{"identity"}\NormalTok{, }\AttributeTok{width =} \DecValTok{1}\NormalTok{) }\SpecialCharTok{+}
  \FunctionTok{coord\_polar}\NormalTok{(}\StringTok{"y"}\NormalTok{, }\AttributeTok{start =} \DecValTok{0}\NormalTok{) }\SpecialCharTok{+} \CommentTok{\# optie voor taartdiagram }
  \FunctionTok{theme\_void}\NormalTok{() }\CommentTok{\# achtergrond weghalen}
\end{Highlighting}
\end{Shaded}

\pandocbounded{\includegraphics[keepaspectratio]{002-02-meetniveaus-visualisaties_files/figure-pdf/unnamed-chunk-4-1.pdf}}

De taartdiagram weergeeft het aandeel per categorie van hoe jongeren
doorgaans aan hun drugs komen. We kunnen de presentatie van taartdiagram
echter nog verbeteren.

\begin{Shaded}
\begin{Highlighting}[]
\CommentTok{\# taartdiagram van hoe drugs worden verkregen}
\FunctionTok{ggplot}\NormalTok{(drugs\_verkrijgen, }\FunctionTok{aes}\NormalTok{(}\AttributeTok{x=} \StringTok{""}\NormalTok{, }\AttributeTok{y =}\NormalTok{ perc, }\AttributeTok{fill =}\NormalTok{ manier\_verkrijgen\_drugs)) }\SpecialCharTok{+} 
  \FunctionTok{geom\_bar}\NormalTok{(}\AttributeTok{stat =} \StringTok{"identity"}\NormalTok{, }\AttributeTok{width =} \DecValTok{1}\NormalTok{, }\AttributeTok{colour =} \StringTok{"white"}\NormalTok{) }\SpecialCharTok{+} \CommentTok{\# witte randen toevoegen}
  \FunctionTok{coord\_polar}\NormalTok{(}\StringTok{"y"}\NormalTok{, }\AttributeTok{start =} \DecValTok{0}\NormalTok{) }\SpecialCharTok{+}
  \FunctionTok{theme\_void}\NormalTok{() }\SpecialCharTok{+} 
  \FunctionTok{geom\_text}\NormalTok{(}\FunctionTok{aes}\NormalTok{(}\AttributeTok{label =}\NormalTok{ perc), }\AttributeTok{position =} \FunctionTok{position\_stack}\NormalTok{(}\AttributeTok{vjust =} \FloatTok{0.5}\NormalTok{)) }\SpecialCharTok{+} \CommentTok{\# percentages als labels toevoegen   }
  \FunctionTok{labs}\NormalTok{(}\AttributeTok{fill =} \StringTok{"Manieren van drugs verkrijgen"}\NormalTok{) }\SpecialCharTok{+} \CommentTok{\# titel voor legenda}
  \FunctionTok{ggtitle}\NormalTok{(}\StringTok{"Taartdiagram met manieren van drugs verkrijgen (\%)"}\NormalTok{) }\CommentTok{\# titel voor figuur}
\end{Highlighting}
\end{Shaded}

\pandocbounded{\includegraphics[keepaspectratio]{002-02-meetniveaus-visualisaties_files/figure-pdf/unnamed-chunk-5-1.pdf}}

We hebben nu een duidelijk inzicht in hoe jongeren die drugs gebruiken
hier aan komen. Zo zien we bijvoorbeeld dat ongeveer xx procent dit via
een dealer koopt, en ongeveer een kwart het koopt via vrienden/kennissen
en een kwart het van deze krijgt.

Dit toont meteen een relevante bevinding uit de literatuur over
drugsgebruik onder jongeren, namelijk dat de helft van de jongeren hun
drugs verkrijgen via ``\emph{sorting}'' (zie
\href{https://journals.sagepub.com/doi/10.1177/003803850203600408}{Parker
et al., 2002}). Dat wil zeggen, dat een bekende uit een vrienden- of
kennissenkring iemand van drugs voorziet eerder dan een dealer.

\subsection{Ordinale variabele}\label{ordinale-variabele}

Een kwalitatieve variabele waarbij er \emph{wel} een bepaalde rangorde
(hiërarchie) tussen de categorieën is, noemen we een \textbf{ordinale}
variabele.

Laten we weer naar een voorbeeld kijken. In de enquete is gevraagd hoe
jongeren hun eigen gezondheid over het algemeen beoordelen. Voor een
ordinale variabele maken we in principe altijd een staafdiagram.

\begin{Shaded}
\begin{Highlighting}[]
\CommentTok{\# staafdiagram van ervaren gezondheid}
\FunctionTok{ggplot}\NormalTok{(data, }\FunctionTok{aes}\NormalTok{(}\AttributeTok{x =}\NormalTok{ ervaren\_gezondheid)) }\SpecialCharTok{+}
  \FunctionTok{geom\_bar}\NormalTok{() }\SpecialCharTok{+}
  \FunctionTok{theme\_bw}\NormalTok{() }\SpecialCharTok{+} 
  \FunctionTok{ylab}\NormalTok{(}\StringTok{"aantal"}\NormalTok{) }\SpecialCharTok{+}
  \FunctionTok{xlab}\NormalTok{ (}\StringTok{"ervaren gezondheid"}\NormalTok{)}
\end{Highlighting}
\end{Shaded}

\pandocbounded{\includegraphics[keepaspectratio]{002-02-meetniveaus-visualisaties_files/figure-pdf/unnamed-chunk-6-1.pdf}}

De staafdiagram toont dat het grootste deel hun gezondheid als ``goed''
beoordeelt. Een relatief klein deel van de jongeren beoordeelt hun
gezondheid als ``matig'' of zelfs ``slecht''.

Dit voorbeeld illustreert een relevant kenmerk van ordinale variabelen:
er is een bepaalde rangorde tussen de categorieen, maar de afstand
ertussen is niet per se gelijk. Zo is de stap van ``matig'' naar
``goed'' bijvoorbeeld kwalitatief anders dan de stap van ``zeer goed''
naar ``uitstekend''. Hierin verschilt een ordinale variabele van een
kwantitatieve variabele, bij de laatste is de afstand tussen waarden
(1-2-3 etc.) in principe altijd gelijk.

\section{Kwantitatieve variabelen}\label{kwantitatieve-variabelen}

Kwantitatieve variabelen zijn numeriek en bestaan dus uit cijfers. De
afstand tussen waarden is altijd gelijk en er zit een rangorde in.
Bijvoorbeeld bij een rapportcijfer van 1 t/m 10 is de afstand tussen 2
en 3 hetzelfde als de afstand tussen 7 en 8. Kwantitatieve variabelen
worden soms ook wel lineaire of continue variabelen genoemd.

\subsection{Interval en ratio
variabelen}\label{interval-en-ratio-variabelen}

Er zijn twee types kwantitatieve variabelen: \textbf{interval} en
\textbf{ratio}. Het verschil tussen deze twee types is dat een interval
variabele geen \emph{absoluut nulpunt} kent en een ratio variabele wel.

Leeftijd is bijvoorbeeld een ratio variabele, je kunt namelijk niet
jonger dan 0 jaar zijn. Geboortejaar, aan de andere kant, is weer een
interval variabele. De christelijke jaartelling begint bij het jaar 0,
maar dat is eigenlijk arbitrair omdat de jaren daarvoor ook geteld
worden. Het specifieke karakter van de variabele bepaalt dus het
meetniveau.

Voor een lineaire variabele met relatief weinig unieke waarden gebruiken
we meestal een staafdiagram. Zie hieronder een voorbeeld van de
variabele leeftijd.

\begin{Shaded}
\begin{Highlighting}[]
\CommentTok{\# staafdiagram van leeftijd}
\FunctionTok{ggplot}\NormalTok{(data, }\FunctionTok{aes}\NormalTok{(}\AttributeTok{x =} \FunctionTok{as.factor}\NormalTok{(leeftijd))) }\SpecialCharTok{+}
  \FunctionTok{geom\_bar}\NormalTok{() }\SpecialCharTok{+}
  \FunctionTok{theme\_bw}\NormalTok{() }\SpecialCharTok{+} 
  \FunctionTok{ylab}\NormalTok{(}\StringTok{"aantal"}\NormalTok{) }\SpecialCharTok{+}
  \FunctionTok{xlab}\NormalTok{ (}\StringTok{"leeftijd"}\NormalTok{)}
\end{Highlighting}
\end{Shaded}

\pandocbounded{\includegraphics[keepaspectratio]{002-02-meetniveaus-visualisaties_files/figure-pdf/unnamed-chunk-7-1.pdf}}

De staafdiagram toont dat relatief veel respondenten tussen 21 en 24
jaar zijn, een leeftijd waarop jongeren veel uitgaan.

\subsection{Histogram}\label{histogram}

Voor een lineaire variabele met juist relatief veel unieke waarden
gebruiken we meestal een histogram. Bij een een histogram worden
bepaalde waardes gegroepeerd tot één staaf, waardoor de presentatie
overzichtelijk blijft. Dit is hieronder weergeven voor de variabele
\emph{gemoed\_dinsdag\_na\_xtc.} Dit is een score tussen de 0 en 100 die
respondenten hebben gegeven hoe zich voelen op de dinsdag na xtc-gebruik
in het weekend, waarbij 0 voor negatief staat en 100 voor positief.

\begin{Shaded}
\begin{Highlighting}[]
\CommentTok{\# histogram van gemoedstoestand op dinsdag na XTC gebruik}
\FunctionTok{ggplot}\NormalTok{(data, }\FunctionTok{aes}\NormalTok{(}\AttributeTok{x =}\NormalTok{ gemoed\_dinsdag\_na\_xtc)) }\SpecialCharTok{+}
  \FunctionTok{geom\_histogram}\NormalTok{(}\AttributeTok{bins =} \DecValTok{50}\NormalTok{) }\SpecialCharTok{+}
  \FunctionTok{theme\_bw}\NormalTok{() }\SpecialCharTok{+} 
  \FunctionTok{ylab}\NormalTok{(}\StringTok{"aantal"}\NormalTok{) }\SpecialCharTok{+}
  \FunctionTok{xlab}\NormalTok{ (}\StringTok{"gemoedstoestand (negatief {-} positief)"}\NormalTok{) }\SpecialCharTok{+}
  \FunctionTok{xlim}\NormalTok{(}\DecValTok{0}\NormalTok{, }\DecValTok{100}\NormalTok{)}
\end{Highlighting}
\end{Shaded}

\pandocbounded{\includegraphics[keepaspectratio]{002-02-meetniveaus-visualisaties_files/figure-pdf/unnamed-chunk-8-1.pdf}}

\section{Samenvatting}\label{samenvatting}

Het meetniveau gaat over het type variabele en wat voor soort informatie
deze bevat. Zoals we later zullen zien, is dit belangrijk om te weten
wat voor soort statistische toets we kunnen uitvoeren. We onderscheiden
de volgende meetniveaus (tussenhaakjes staat welke figuur we meestal
gebruiken):

\begin{itemize}
\item
  nominaal (staaf- of taartdiagram)
\item
  ordinaal (staafdiagram)
\item
  interval (staafdiagram of histogram)
\item
  ratio (staafdiagram of histogram)
\end{itemize}

\chapter{Centrum- en spreidingsmaten}\label{centrum--en-spreidingsmaten}

Inmiddels zijn we bekend met datasets, observaties, variabelen,
meetniveaus en figuren. We hebben ook benoemd dat statistiek draait om
het analyseren van \emph{variatie} met behulp van de variabelen die in
onze dataset zitten. Denk aan: variatie in drugsgebruik, leeftijden en
opleidingsniveau. Om de variatie te analyseren maken we gebruik van
diverse centrum- en spreidingsmaten. De meeste hiervan zul je herkennen
van eerdere wiskundelessen die je hebt gehad.

\section{Centrummaten}\label{centrummaten}

\subsection{Gemiddelde}\label{gemiddelde}

Het gemiddelde van een variabele berekenen we simpelweg door alle
waardes bij elkaar op te tellen en dan de som te delen door het aantal
observaties \emph{n}. Dit gaat middels de volgende formule.

\[
\bar{x}=\frac{\sum_{i=1}^{n}{x_i}}{n}
\]

\[
S^2=\frac{\sum_{i=1}^{n}{(x_i-\bar{x})^2}}{n-1}
\]

\[
\sigma^2=\frac{\sum_{i=1}^{n}(x_i-\mu)^2}{N}
\]

\[
z=\frac{x-\bar{x}}{S}
\]

\[
s.e. =\frac{S}{\sqrt{n}}
\]

Formules berekenen en statistische toetsen uitvoeren kunnen we ook
handmatig in R doen. Dit is leerzaam, omdat het intuïtie creëert voor
hoe je statististisch onderzoek moet uitvoeren. In onderstaand voorbeeld
berekenen we in R het gemiddelde van een cijferreeks.

\begin{Shaded}
\begin{Highlighting}[]
\NormalTok{(}\DecValTok{3} \SpecialCharTok{+} \DecValTok{5} \SpecialCharTok{+} \DecValTok{5} \SpecialCharTok{+} \DecValTok{7} \SpecialCharTok{+} \DecValTok{10}\NormalTok{) }\SpecialCharTok{/} \DecValTok{5}
\end{Highlighting}
\end{Shaded}

\begin{verbatim}
[1] 6
\end{verbatim}

Het gemiddelde van deze cijferreeks met 5 getallen is dus 6. In plaats
van handmatig kun je het ook door een functie in R doen:

\begin{Shaded}
\begin{Highlighting}[]
\CommentTok{\# vector aanmaken}
\NormalTok{x }\OtherTok{\textless{}{-}} \FunctionTok{c}\NormalTok{(}\DecValTok{3}\NormalTok{, }\DecValTok{5}\NormalTok{, }\DecValTok{5}\NormalTok{, }\DecValTok{7}\NormalTok{, }\DecValTok{10}\NormalTok{)}
\FunctionTok{print}\NormalTok{(x)}
\end{Highlighting}
\end{Shaded}

\begin{verbatim}
[1]  3  5  5  7 10
\end{verbatim}

\begin{Shaded}
\begin{Highlighting}[]
\CommentTok{\# gemiddelde berekenen}
\FunctionTok{mean}\NormalTok{(x)}
\end{Highlighting}
\end{Shaded}

\begin{verbatim}
[1] 6
\end{verbatim}

Eerst hebben we een vector, of eigenlijk een variabele, aangemaakt met
de vijf waardes. Via de functie \emph{mean()} bereken je het gemiddelde.

\subsection{Mediaan}\label{mediaan}

De mediaan is het getal waar het midden van het aantal observaties is
bereikt (50\%), oftewel de middelste observatie.

Bij de cijferreeks (3, 5, \ul{5}, 7, 10) is dit 5.

Bij een even aantal cijfers zoals (3, 5, 5, 6, 7, 10) neem je het midden
van de middelste twee getallen. In dit geval is dat 5,5 (tussen 5 en 6).

In R bereken je de mediaan via de functie \emph{median().}

\begin{Shaded}
\begin{Highlighting}[]
\CommentTok{\# oneven aantal observaties}
\FunctionTok{median}\NormalTok{(x)}
\end{Highlighting}
\end{Shaded}

\begin{verbatim}
[1] 5
\end{verbatim}

\begin{Shaded}
\begin{Highlighting}[]
\CommentTok{\# even aantal observaties}
\NormalTok{y }\OtherTok{\textless{}{-}} \FunctionTok{c}\NormalTok{(}\DecValTok{3}\NormalTok{, }\DecValTok{5}\NormalTok{, }\DecValTok{5}\NormalTok{, }\DecValTok{6}\NormalTok{, }\DecValTok{7}\NormalTok{, }\DecValTok{10}\NormalTok{)}
\FunctionTok{print}\NormalTok{(y)}
\end{Highlighting}
\end{Shaded}

\begin{verbatim}
[1]  3  5  5  6  7 10
\end{verbatim}

\begin{Shaded}
\begin{Highlighting}[]
\FunctionTok{median}\NormalTok{(y)}
\end{Highlighting}
\end{Shaded}

\begin{verbatim}
[1] 5.5
\end{verbatim}

\subsection{Modus}\label{modus}

De modus is de waarde of categorie die het vaakst voorkomt in een
variabele.

Bij de cijferreeks (3, 5, 5, 7, 10) is dit 5 omdat deze tweemaal
voorkomt en de rest slechts een keer.

R bereken je de modus via de functie \emph{mfv().}

\begin{Shaded}
\begin{Highlighting}[]
\CommentTok{\# package laden}
\FunctionTok{library}\NormalTok{(modeest)}

\CommentTok{\# modus berekenen}
\FunctionTok{mfv}\NormalTok{(x)}
\end{Highlighting}
\end{Shaded}

\begin{verbatim}
[1] 5
\end{verbatim}

\part{Kans en steekproeven}

\chapter{Steekproeven}\label{steekproeven}

\chapter{Kans}\label{kans}

\part{Hypotheses}

\part{Samenhang tussen variabelen}

\part{Colleges}

\chapter{Collegeslides}\label{sec-colleges}

Week 1: \href{collegeslides/college_1.html}{College 1}

\part{Opdrachten}

\chapter{Opdrachten}\label{opdrachten-1}

Deze week gaan jullie aan de slag met \ldots{}

\cleardoublepage
\phantomsection
\addcontentsline{toc}{part}{Appendices}
\appendix

\chapter{Symbolen en formules}\label{symbolen-en-formules}

\section{Belangrijke symbolen en
formules}\label{belangrijke-symbolen-en-formules}

Hier vind je een overzicht van alle belangrijke symbolen en formules die
we gebruiken.

{\def\LTcaptype{none} % do not increment counter
\begin{longtable}[]{@{}
  >{\raggedright\arraybackslash}p{(\linewidth - 2\tabcolsep) * \real{0.1667}}
  >{\raggedright\arraybackslash}p{(\linewidth - 2\tabcolsep) * \real{0.8333}}@{}}
\toprule\noalign{}
\begin{minipage}[b]{\linewidth}\raggedright
Symbool
\end{minipage} & \begin{minipage}[b]{\linewidth}\raggedright
Betekenis
\end{minipage} \\
\midrule\noalign{}
\endhead
\bottomrule\noalign{}
\endlastfoot
\(\sum_{i=1}^n\) & Griekse hoofdletter sigma, staat voor `sommatie':
alles bij elkaar optellen. De \(i\) betekent dat je begint met tellen
bij de eerste waarneming, de \(n\) dat je stopt bij de laatste. \\
\(\alpha\) & Alpha, het significantieniveau \\
\(\mu\) & Mu, het gemiddelde in de \emph{populatie} (vaak onbekend) \\
\(\sigma\) & Sigma, de standaarddeviatie in een \emph{populatie} (vaak
onbekend) \\
\(\sigma^2\) & Sigma kwadraat, de variantie in een \emph{populatie} \\
\(\sigma_{\bar{x}}\) & Sigma mean, standaardfout vna het gemiddelde \\
\end{longtable}
}




\end{document}
